% HQIV_LEAN: factor-search drivers — algorithm description and formalized proofs
% Build: pdflatex (twice for references) or latexmk
\documentclass[11pt,a4paper]{article}

\usepackage{amsmath,amssymb,amsthm,mathtools}
\usepackage{microtype}
\usepackage[hidelinks]{hyperref}
\usepackage[margin=1in]{geometry}
\usepackage{enumitem}
\usepackage{url}

\newcommand{\N}{\mathbb{N}}
\newcommand{\Q}{\mathbb{Q}}
\newcommand{\R}{\mathbb{R}}
\newcommand{\dist}{\operatorname{dist}}
\newcommand{\bitlen}{\operatorname{bitlen}}

\theoremstyle{plain}
\newtheorem{theorem}{Theorem}[section]
\newtheorem{lemma}[theorem]{Lemma}
\newtheorem{proposition}[theorem]{Proposition}

\theoremstyle{definition}
\newtheorem{definition}[theorem]{Definition}

\theoremstyle{remark}
\newtheorem{remark}[theorem]{Remark}

\title{Factor Search in HQIV\textunderscore LEAN:\\
Algorithms, Combinatorial Certificates, and Formalized Proofs}
\author{HQIV\_LEAN repository (machine-checked fragments)}
\date{\today}

\begin{document}
\maketitle

\begin{abstract}
This note records the \emph{algorithmic} structure of two HQIV factor-search drivers
(``integrated'' sparse OSH and ``monolithic'' geometric factorizer) and summarizes
what is actually proved in the Lean~4 library \texttt{HQIV\_LEAN}: sound divisibility
certificates, gcd/remainder equivalence, sparse-register list algebra, explicit
per-step charging schemes, and a polynomial \emph{template} bound for a combined
integrated cost model. We explicitly separate these results from open claims:
no theorem here establishes that either driver finds a nontrivial factor on every
composite input, nor that wall-clock time scales as a particular function of $\log n$
without additional machine hypotheses.
\end{abstract}

% Shared reader contract: patch QFT + discrete numerics.
% From papers/*.tex:     % Shared reader contract: patch QFT + discrete numerics.
% From papers/*.tex:     % Shared reader contract: patch QFT + discrete numerics.
% From papers/*.tex:     \input{include/patch_theory_messaging.tex}
% From papers/paper/*.tex: \input{../include/patch_theory_messaging.tex}

\paragraph{Patch QFT and discrete numerics (reader contract).}
HQIV is formulated as a \emph{patch theory}: quantum fields, actions, and physical readouts are attached to discrete patches---shell indices $m\in\mathbb{N}$, finite spacetime index types (e.g.\ $\mathrm{Fin}\,4$), and horizon-limited charts---rather than to a hypothetical fundamental smooth spacetime obtained as a mesh-refinement or ultraviolet limit.
Accordingly, when this text refers to ``QFT,'' it means \emph{patch QFT}: patching rules, finite measurement ledgers, and variational identities on the discrete spine; it is not the claim that the theory is \emph{defined} by recovering ordinary continuum quantum field theory through such a limit.
The numbers that admit the sharpest statements---mass and coupling ladders, curvature ratios, shell thermodynamics, and rational witnesses in \texttt{HQIV\_LEAN}---are objects of \emph{discrete mathematics} on $\mathbb{N}$ (combinatorial growth, cumulative simplex ledgers) together with the ordered field $\mathbb{R}$ as used in the proof assistant for analysis, bounds, and phenomenological display.
Smooth-manifold or continuum field notation, when it appears, is a \emph{translation layer} for comparison with standard physics literature; it does not supersede the discrete definitions.

% From papers/paper/*.tex: % Shared reader contract: patch QFT + discrete numerics.
% From papers/*.tex:     \input{include/patch_theory_messaging.tex}
% From papers/paper/*.tex: \input{../include/patch_theory_messaging.tex}

\paragraph{Patch QFT and discrete numerics (reader contract).}
HQIV is formulated as a \emph{patch theory}: quantum fields, actions, and physical readouts are attached to discrete patches---shell indices $m\in\mathbb{N}$, finite spacetime index types (e.g.\ $\mathrm{Fin}\,4$), and horizon-limited charts---rather than to a hypothetical fundamental smooth spacetime obtained as a mesh-refinement or ultraviolet limit.
Accordingly, when this text refers to ``QFT,'' it means \emph{patch QFT}: patching rules, finite measurement ledgers, and variational identities on the discrete spine; it is not the claim that the theory is \emph{defined} by recovering ordinary continuum quantum field theory through such a limit.
The numbers that admit the sharpest statements---mass and coupling ladders, curvature ratios, shell thermodynamics, and rational witnesses in \texttt{HQIV\_LEAN}---are objects of \emph{discrete mathematics} on $\mathbb{N}$ (combinatorial growth, cumulative simplex ledgers) together with the ordered field $\mathbb{R}$ as used in the proof assistant for analysis, bounds, and phenomenological display.
Smooth-manifold or continuum field notation, when it appears, is a \emph{translation layer} for comparison with standard physics literature; it does not supersede the discrete definitions.


\paragraph{Patch QFT and discrete numerics (reader contract).}
HQIV is formulated as a \emph{patch theory}: quantum fields, actions, and physical readouts are attached to discrete patches---shell indices $m\in\mathbb{N}$, finite spacetime index types (e.g.\ $\mathrm{Fin}\,4$), and horizon-limited charts---rather than to a hypothetical fundamental smooth spacetime obtained as a mesh-refinement or ultraviolet limit.
Accordingly, when this text refers to ``QFT,'' it means \emph{patch QFT}: patching rules, finite measurement ledgers, and variational identities on the discrete spine; it is not the claim that the theory is \emph{defined} by recovering ordinary continuum quantum field theory through such a limit.
The numbers that admit the sharpest statements---mass and coupling ladders, curvature ratios, shell thermodynamics, and rational witnesses in \texttt{HQIV\_LEAN}---are objects of \emph{discrete mathematics} on $\mathbb{N}$ (combinatorial growth, cumulative simplex ledgers) together with the ordered field $\mathbb{R}$ as used in the proof assistant for analysis, bounds, and phenomenological display.
Smooth-manifold or continuum field notation, when it appears, is a \emph{translation layer} for comparison with standard physics literature; it does not supersede the discrete definitions.

% From papers/paper/*.tex: % Shared reader contract: patch QFT + discrete numerics.
% From papers/*.tex:     % Shared reader contract: patch QFT + discrete numerics.
% From papers/*.tex:     \input{include/patch_theory_messaging.tex}
% From papers/paper/*.tex: \input{../include/patch_theory_messaging.tex}

\paragraph{Patch QFT and discrete numerics (reader contract).}
HQIV is formulated as a \emph{patch theory}: quantum fields, actions, and physical readouts are attached to discrete patches---shell indices $m\in\mathbb{N}$, finite spacetime index types (e.g.\ $\mathrm{Fin}\,4$), and horizon-limited charts---rather than to a hypothetical fundamental smooth spacetime obtained as a mesh-refinement or ultraviolet limit.
Accordingly, when this text refers to ``QFT,'' it means \emph{patch QFT}: patching rules, finite measurement ledgers, and variational identities on the discrete spine; it is not the claim that the theory is \emph{defined} by recovering ordinary continuum quantum field theory through such a limit.
The numbers that admit the sharpest statements---mass and coupling ladders, curvature ratios, shell thermodynamics, and rational witnesses in \texttt{HQIV\_LEAN}---are objects of \emph{discrete mathematics} on $\mathbb{N}$ (combinatorial growth, cumulative simplex ledgers) together with the ordered field $\mathbb{R}$ as used in the proof assistant for analysis, bounds, and phenomenological display.
Smooth-manifold or continuum field notation, when it appears, is a \emph{translation layer} for comparison with standard physics literature; it does not supersede the discrete definitions.

% From papers/paper/*.tex: % Shared reader contract: patch QFT + discrete numerics.
% From papers/*.tex:     \input{include/patch_theory_messaging.tex}
% From papers/paper/*.tex: \input{../include/patch_theory_messaging.tex}

\paragraph{Patch QFT and discrete numerics (reader contract).}
HQIV is formulated as a \emph{patch theory}: quantum fields, actions, and physical readouts are attached to discrete patches---shell indices $m\in\mathbb{N}$, finite spacetime index types (e.g.\ $\mathrm{Fin}\,4$), and horizon-limited charts---rather than to a hypothetical fundamental smooth spacetime obtained as a mesh-refinement or ultraviolet limit.
Accordingly, when this text refers to ``QFT,'' it means \emph{patch QFT}: patching rules, finite measurement ledgers, and variational identities on the discrete spine; it is not the claim that the theory is \emph{defined} by recovering ordinary continuum quantum field theory through such a limit.
The numbers that admit the sharpest statements---mass and coupling ladders, curvature ratios, shell thermodynamics, and rational witnesses in \texttt{HQIV\_LEAN}---are objects of \emph{discrete mathematics} on $\mathbb{N}$ (combinatorial growth, cumulative simplex ledgers) together with the ordered field $\mathbb{R}$ as used in the proof assistant for analysis, bounds, and phenomenological display.
Smooth-manifold or continuum field notation, when it appears, is a \emph{translation layer} for comparison with standard physics literature; it does not supersede the discrete definitions.


\paragraph{Patch QFT and discrete numerics (reader contract).}
HQIV is formulated as a \emph{patch theory}: quantum fields, actions, and physical readouts are attached to discrete patches---shell indices $m\in\mathbb{N}$, finite spacetime index types (e.g.\ $\mathrm{Fin}\,4$), and horizon-limited charts---rather than to a hypothetical fundamental smooth spacetime obtained as a mesh-refinement or ultraviolet limit.
Accordingly, when this text refers to ``QFT,'' it means \emph{patch QFT}: patching rules, finite measurement ledgers, and variational identities on the discrete spine; it is not the claim that the theory is \emph{defined} by recovering ordinary continuum quantum field theory through such a limit.
The numbers that admit the sharpest statements---mass and coupling ladders, curvature ratios, shell thermodynamics, and rational witnesses in \texttt{HQIV\_LEAN}---are objects of \emph{discrete mathematics} on $\mathbb{N}$ (combinatorial growth, cumulative simplex ledgers) together with the ordered field $\mathbb{R}$ as used in the proof assistant for analysis, bounds, and phenomenological display.
Smooth-manifold or continuum field notation, when it appears, is a \emph{translation layer} for comparison with standard physics literature; it does not supersede the discrete definitions.


\paragraph{Patch QFT and discrete numerics (reader contract).}
HQIV is formulated as a \emph{patch theory}: quantum fields, actions, and physical readouts are attached to discrete patches---shell indices $m\in\mathbb{N}$, finite spacetime index types (e.g.\ $\mathrm{Fin}\,4$), and horizon-limited charts---rather than to a hypothetical fundamental smooth spacetime obtained as a mesh-refinement or ultraviolet limit.
Accordingly, when this text refers to ``QFT,'' it means \emph{patch QFT}: patching rules, finite measurement ledgers, and variational identities on the discrete spine; it is not the claim that the theory is \emph{defined} by recovering ordinary continuum quantum field theory through such a limit.
The numbers that admit the sharpest statements---mass and coupling ladders, curvature ratios, shell thermodynamics, and rational witnesses in \texttt{HQIV\_LEAN}---are objects of \emph{discrete mathematics} on $\mathbb{N}$ (combinatorial growth, cumulative simplex ledgers) together with the ordered field $\mathbb{R}$ as used in the proof assistant for analysis, bounds, and phenomenological display.
Smooth-manifold or continuum field notation, when it appears, is a \emph{translation layer} for comparison with standard physics literature; it does not supersede the discrete definitions.


\tableofcontents

\section{Introduction}

The repository implements classical search loops paired with Lean mirrors. The
mathematical content that is \emph{proved} falls into three classes:
\begin{enumerate}[label=(\roman*)]
\item \textbf{Sound certificates:} once an integer passes a modular test, the
      factorization identities recorded in Lean follow from standard arithmetic.
\item \textbf{Combinatorial bookkeeping:} sparse simulation steps expand and prune
      lists; Lean bounds list lengths and sums of lengths.
\item \textbf{Cost scaffolding:} under explicit definitions of ``fuel'' and per-step
      charges, Lean proves inequalities such as a degree-$6$ polynomial bound in
      $\lfloor \log_2(n+1)\rfloor + 1$.
\end{enumerate}

What is \emph{not} proved (unless explicitly stated as a conjecture or roadmap item)
includes completeness of search, correctness of floating-point angle rounding, and
primality of factors returned by probabilistic engines.

\section{Algorithms (informal)}

\subsection{Integrated driver (\texttt{hqiv\_osh\_integrated\_driver.py})}

At a high level the script:
\begin{itemize}
\item Peels factors of two from the input, working on an odd core.
\item Builds a sparse harmonic register (cutoff $L$, basis size $(L{+}1)^2$).
\item Runs sparse gate steps: expand support, apply a digital gate, detect flipped
      indices, prune to flipped support (the ``sparse Shor-style'' spine).
\item Maintains phase / angle channels that map indices to integer candidates; on a
      hit, checks classical divisibility.
\end{itemize}

The Lean module \texttt{HQIVOSHIntegratedFactorDriver} names the discrete maps
(\texttt{rootCap}, \texttt{idxAngleCandidate}, \texttt{wrapIdx}, ket fallback) and
packages divisor witnesses after peeling.

\subsection{Monolithic geometric factorizer (\texttt{monolithic\_geometric\_factorizer.py})}

The monolithic script:
\begin{itemize}
\item Uses dynamic decimal precision tied to the bit length of the odd core.
\item Maps rational turns $a/n$ to integer candidates via a root cap
      $\max(2,\lfloor\sqrt{n}\rfloor)$.
\item Applies multiplication gates with inverse strength keyed off register bit sizes
      (``lowest register'' variant in Lean).
\item Prunes candidate lines by an age threshold derived from
      $\bitlen(|c-n|)+1$.
\item Caps the number of parallel registers at cube-root scale
      $\min(128,\lfloor n^{1/3}\rfloor)$ (after peeling twos).
\item Accepts a factor only on a purely integer test: $1<d<\text{odd}$ and
      $\text{odd}\bmod d = 0$.
\end{itemize}

The Lean file \texttt{MonolithicGeometricFactorizer} mirrors these discrete parameters
and isolates real-number drift resources (\texttt{MonolithicFracResources}) as a
\emph{template} for aligning floating-point $\theta/(2\pi)$ with floor slots.

\section{Divisibility certificates and GCD}

\subsection{Odd-core witness}

\begin{definition}[Odd-core factor witness]
Given odd $n\in\N$, an \emph{odd-core factor witness} consists of $d\in\N$ with
$1<d<n$ and $d\mid n$.
\end{definition}

Lean encodes this as \texttt{OddCoreFactorWitness} (\texttt{HQIVOSHIntegratedFactorDriver}).
Reconstruction $n = d\cdot (n/d)$ is \texttt{OddCoreFactorWitness.reconstructs}.

\subsection{Peeling powers of two}

If $\mathrm{orig} = 2^k \cdot \mathrm{odd}$ and $d\mid \mathrm{odd}$, then
$2^k d \mid \mathrm{orig}$ (\texttt{two\_pow\_mul\_dvd\_of\_dvd\_odd}).

\subsection{GCD and division (\texttt{FactorDivisibilityBridge})}

\begin{proposition}[Division and quotient]
For all $d,n\in\N$, $d\mid n$ if and only if $n = d\cdot (n/d)$.
\end{proposition}

\begin{proposition}[GCD as the same certificate class]
If $1 < \gcd(x,\mathrm{odd}) < \mathrm{odd}$, then $\gcd(x,\mathrm{odd})$ yields an
odd-core witness. If already $d\mid \mathrm{odd}$, then $\gcd(d,\mathrm{odd})=d$.
\end{proposition}

\begin{remark}
These statements are pure arithmetic; they do \emph{not} guarantee that a particular
random $x$ produces a nontrivial gcd. A concrete counterexample is
$\gcd(2,15)=1$ while $3\mid 15$ (\texttt{gcd\_probe\_can\_miss\_nontrivial\_factor}).
\end{remark}

\subsection{Monolithic acceptance}

\begin{definition}
\texttt{monolithicAcceptsDivisor}$(n,d)$ $\equiv$ $1<d<n$ and $n\bmod d = 0$.
\end{definition}

\begin{proposition}
\texttt{monolithicAcceptsDivisor}$(n,d)$ $\Leftrightarrow$ $1<d<n$ and $d\mid n$
(\texttt{monolithicAcceptsDivisor\_iff\_dvd\_and\_bounds}).
\end{proposition}

\begin{proposition}[Compositeness]
If $1<n$ and \texttt{monolithicAcceptsDivisor}$(n,d)$, then $n$ is not prime
(\texttt{not\_prime\_of\_monolithicAccepts}).
\end{proposition}

\section{Sparse OSH oracle: list algebra}

Let a \texttt{SparseRegister} be a finite list of indexed amplitudes. One step:
\begin{enumerate}
\item \texttt{causalExpandSupport} doubles support in length (horizon-causal expansion).
\item \texttt{applyGateSparse} maps to length $2|r|$ (\texttt{applyGateSparse\_length\_eq\_two\_mul}).
\item \texttt{detectFlippedKets} lists indices that changed between ``before'' and ``after''.
\item \texttt{pruneToFlipped} filters to a flipped index set; length never increases
      (\texttt{pruneToFlipped\_length\_le}).
\end{enumerate}

\begin{proposition}[Integrated step length]
After evolve then prune,
$|\mathrm{prune}| \le 2 \cdot |r|$ (\texttt{integrated\_pruned\_sparse\_step\_length\_le}).
\end{proposition}

\begin{proposition}[Triple sum bound]
With $e=\texttt{applyGateSparse}(g,r)$ and $f=\texttt{detectFlippedKets}(r,e)$,
\[
|r| + |e| + |f| \le 6|r|
\]
(\texttt{sparse\_step\_list\_sum\_le\_six\_mul} in \texttt{FactorSearchCostModel}).
\end{proposition}

\section{Integrated driver maps and bounds}

\subsection{Root cap and angle candidate}

$\texttt{rootCap}(n) = \max(2,\lfloor\sqrt{n}\rfloor)$, always $\ge 2$.

The discrete angle candidate (for basis $b$ and index $\texttt{idx}$) satisfies
\[
\texttt{idxAngleCandidate}(n,b,\texttt{idx}) \le 2 + \texttt{idx}\cdot \texttt{rootCap}(n)
\]
(\texttt{idxAngleCandidate\_le\_two\_add\_idx\_mul\_rootCap}).

\subsection{Wrapped indices}

$\texttt{wrapIdx}(L,\texttt{idx}) < (L+1)^2 = \texttt{sparseBasisCard}(L)$.

\subsection{Ket fallback alignment}

The integrated ket fallback agrees with the gate-frontier cofactor map on wrapped
indices (\texttt{ketLinearFallback\_eq\_cofactorCandidate}), and when
$2 \le \texttt{qSpan}(n)$, candidates stay inside the $Q$-shell bound
(\texttt{ketLinearFallbackCandidate\_le\_qSpan}).

\section{Cost model and polynomial template}

\subsection{Bit-scale parameter}

$\texttt{inputLog2}(n) := \lfloor \log_2(n+1)\rfloor$ and
$d := \texttt{inputLogDim}(n) = \texttt{inputLog2}(n)+1$.

\subsection{Default harmonic cutoff}

$\texttt{defaultHarmonicCutoff}(n) = \max(4,\lfloor\sqrt{n}\rfloor)$, hence
$L+1 \le 5+\lfloor\sqrt{n}\rfloor$ and basis cardinality $(L+1)^2 \le (5+\lfloor\sqrt{n}\rfloor)^2$
for this default.

\subsection{Combined per-step cap}

The file \texttt{FactorSearchCostModel} defines a conservative combined cap
\texttt{integratedCombinedPerStepCap} that sums:
\begin{itemize}
\item a sparse list charge modeled as $6\cdot d^3$ in terms of $d=\texttt{inputLogDim}(n)$;
\item a phase-channel charge with history length bounded by fuel and extract cost $f^2$.
\end{itemize}

\begin{theorem}[Per-step domination]
$\texttt{integratedCombinedPerStepCap}(n) \le 12 \cdot d^4$
(\texttt{integratedCombinedPerStepCap\_le\_twelve\_mul\_inputLogDim\_pow\_four}).
\end{theorem}

\subsection{Total work under polylog fuel}

Let $\texttt{fuelPolyLogSq}(n) = d^2$. If $\texttt{fuel} \le \texttt{fuelPolyLogSq}(n)$,
then naive total work charged at \texttt{integratedCombinedPerStepCap} satisfies
\[
\texttt{fuelTotalWork}(\texttt{fuel},\texttt{integratedCombinedPerStepCap}(n))
\;\le\; 12 \cdot d^6
\]
(\texttt{total\_integrated\_work\_le\_twelve\_mul\_inputLogDim\_pow\_six}).

\begin{remark}
This is a \emph{rigorous inequality in the cost model} once fuel and charges are
instantiated as above. It is not, by itself, a wall-clock theorem: real implementations
allocate dense intermediates, use floating-point, and may terminate earlier or later
than the fuel model.
\end{remark}

\section{Monolithic discrete policy (selected lemmas)}

\subsection{Cube-root floor}

$\texttt{floorCbrt}(n)$ is the greatest $r$ with $r^3 \le n$; it satisfies the
standard maximality properties (\texttt{floorCbrt\_cube\_le}, \texttt{floorCbrt\_lt\_succ\_cube}).

\subsection{Dynamic precision}

$\texttt{dynamicPrecisionDigits}(n) \le 2\cdot \texttt{inputLog2}(n) + 11$
(\texttt{dynamicPrecisionDigits\_le\_two\_mul\_inputLog\_add\_eleven}).

\subsection{Prune age threshold}

For candidate $c$ and target $n$, write $d=\dist(c,n)$ and define bit-length of $d$ as
$0$ if $d=0$ else $\lfloor\log_2 d\rfloor+1$. Then
$\texttt{pruneAgeThreshold}(n,c)$ is this bit-length plus $1$.

For $c\le n$, $\texttt{pruneAgeThreshold}(n,c) \le \texttt{inputLog2}(n)+3$
(\texttt{pruneAgeThreshold\_le\_inputLog\_add\_three}).

\subsection{Mod-periodicity certificate}

If $f(t+p)\equiv f(t) \pmod m$ for all $t$, then residues at ages reduce to
$f(\mathrm{age}\bmod p)$ (\texttt{mod\_residue\_eq\_mod\_age}) and equal ages with the
same remainder mod $p$ yield equal residues (\texttt{age\_redundant\_mod\_obs\_of\_global\_period}).

\section{Counterexamples (boundary effects)}

Concrete arithmetic facts proved with \texttt{native\_decide}:
\begin{itemize}
\item Distances $\dist(993,1000)=7$ and $\dist(992,1000)=8$ but different
      $\texttt{pruneAgeThreshold}$ values $4$ vs.\ $5$
      (\texttt{pruneAgeThreshold\_jumps\_across\_dist\_boundary}).
\item Register bit-size proxies differ across $7$ vs.\ $8$
      (\texttt{registerBitSize\_jumps\_at\_eight}).
\item Inverse lowest-register strength differs for pairs $(7,100)$ vs.\ $(8,100)$
      (\texttt{strengthInvLowestRegister\_jumps\_when\_low\_coordinate\_crosses\_eight}).
\end{itemize}

\section{$G_2$-cubic factor readout for Legendre-anomalous primes}
\label{sec:g2-cubic-factor-readout}

This section records a closed-form identity discovered while probing the
\texttt{so8\_factor.py} prototype. It is \emph{not yet} formalized in Lean but is
fully verified numerically (see \texttt{tests/test\_factorization.py::test\_g2\_cubic\_invariant\_*}).

\paragraph{Setup.}
Let $N = p \cdot q$ with $p < q$ both odd primes. Choose Lagrange four-square
representations
\[
p = a_0^2 + a_1^2 + a_2^2 + a_3^2,\qquad
q = b_0^2 + b_1^2 + b_2^2 + b_3^2,
\]
and form octonions
\[
\alpha = a_0 + a_1 e_1 + a_2 e_2 + a_3 e_3 \in \mathbb{H},\qquad
\gamma = (b_0 + b_1 e_1 + b_2 e_2 + b_3 e_3) \cdot e_4 \in \mathbb{H}\!\cdot\!e_4,
\]
so that $\mathcal{N}(\alpha) = p,\ \mathcal{N}(\gamma) = q,\ \mathcal{N}(\alpha\!\cdot\!\gamma) = N$.
Let $\Phi(x,y,z) = \langle x,\, \mathrm{Im}(y\!\cdot\!z)\rangle$ denote the
$G_2$-invariant trilinear cubic form on $\operatorname{Im}(\mathbb{O})$.

\begin{theorem}[$G_2$-cubic factor readout]\label{thm:g2-cubic}
With notation as above, the Cayley--Dickson identity $\gamma\!\cdot\!\alpha = (\beta\!\cdot\!\bar\alpha)\cdot e_4$
implies the closed-form identities
\begin{align}
\Phi(\gamma,\gamma,\alpha) &= q \cdot a_0, \label{eq:g2-cubic}\\
\|[\alpha,\gamma]\|^2 &= 4\,q\,(p - a_0^2). \label{eq:comm-norm}
\end{align}
In particular $\Phi(\alpha,\alpha,\gamma) = 0$ because both $\alpha\!\cdot\!\gamma$ and
$\gamma\!\cdot\!\alpha$ land in the coassociative block $\operatorname{span}(e_4,\ldots,e_7)$,
orthogonal to $\alpha \in \mathbb{H}$.
\end{theorem}

\begin{proof}[Sketch]
The Cayley--Dickson rule on $\mathbb{O} = \mathbb{H} \oplus \mathbb{H}\!\cdot\!e_4$ gives
$\gamma\!\cdot\!\alpha = (\beta\!\cdot\!\bar\alpha)\!\cdot\!e_4$.
Right-multiplication by $e_4$ is an isometry of the coassociative block, so
$\langle \beta\!\cdot\!e_4,\, (\beta\!\cdot\!\bar\alpha)\!\cdot\!e_4\rangle
 = \langle \beta,\, \beta\!\cdot\!\bar\alpha\rangle = \mathrm{Re}(\bar\beta\!\cdot\!\beta\!\cdot\!\bar\alpha)
 = \mathcal{N}(\beta) \cdot \mathrm{Re}(\bar\alpha) = q\, a_0$, proving \eqref{eq:g2-cubic}.
Equation \eqref{eq:comm-norm} follows from direct expansion of
$\alpha\!\cdot\!\gamma - \gamma\!\cdot\!\alpha$ in the coassociative block.
\end{proof}

\paragraph{Legendre's three-square obstruction.}
By Legendre's three-square theorem, $p$ is a sum of three squares iff
$p \not\equiv 7 \pmod 8$. Thus, for primes $p \equiv 7 \pmod 8$ every
four-square representation has $a_0 \ge 1$, and the readout \eqref{eq:g2-cubic}
is \emph{forced} nontrivial:
\[
p \equiv 7 \pmod 8 \implies a_0 \neq 0 \implies q = \Phi(\gamma,\gamma,\alpha)/a_0.
\]
This recovers the cofactor $q$ without classical division, given $p$.

\paragraph{Forbidden cubic addition mod 7.}
The cubes modulo $7$ are $\{0,1,6\}$, hence sums of two cubes modulo $7$ lie in
$\{0,1,2,5,6\}$; the values $\{3,4\}$ are \emph{forbidden}. Substituting
$\Phi(\gamma,\gamma,\alpha) = q\,a_0$ shows that, when $a_0 = 1$,
\[
\Phi(\gamma,\gamma,\alpha) \bmod 7 \in \{3, 4\}\iff q \bmod 7 \in \{3, 4\}.
\]
This is a sieve criterion: among the semiprimes for which Legendre triggers,
the cofactor's mod-7 class is read off from the $G_2$-cubic invariant directly.

\paragraph{Empirically verified cases ($p < q$, $p \equiv 7 \pmod 8$ or
$p$ chosen with $a_0 = 1$).}
\begin{center}
\begin{tabular}{rrrrr}
\hline
$N$ & $p$ & $q$ & $\Phi(\gamma,\gamma,\alpha)$ & $q \bmod 7$ \\ \hline
$77$ & $7$ & $11$ & $11$ & $4$ (forbidden) \\
$91$ & $7$ & $13$ & $13$ & $6$ \\
$221$ & $13$ & $17$ & $17$ & $3$ (forbidden) \\
$247$ & $13$ & $19$ & $19$ & $5$ \\
$667$ & $23$ & $29$ & $29$ & $1$ \\ \hline
\end{tabular}
\end{center}

\paragraph{Scope.}
Theorem~\ref{thm:g2-cubic} does \emph{not} produce a polynomial-time factoring
algorithm: computing $\Phi(\gamma,\gamma,\alpha)$ requires constructing $\alpha,
\gamma$, which presupposes the four-square representations of $p$ and $q$. The
result is structural: it pins the cofactor of a Legendre-anomalous semiprime
to a single closed-form readout off the $G_2$-invariant cubic, and ties the
mod-7 ``forbidden cubic addition'' classification to the residue class of $q$.
This is the analytic underpinning we use to justify the $\mathrm{so}(8)+\Delta$
factor layer in \texttt{factorization.py}; a Lean formalization of
\eqref{eq:g2-cubic} and \eqref{eq:comm-norm} is a natural target for the next
formalization pass.

\subsection{Three-layer structural picture: $\mathrm{so}(7) \subset \mathrm{so}(8)$ and the
harmonic factor basis}
\label{sec:harmonic-factor-basis}

The cubic readout admits an additional layer of geometric symmetry that exposes
all four Lagrange coefficients of $\alpha$ via three harmonic amplitudes.

\paragraph{Layer 1 ($\mathrm{so}(7)$-stabilizer of $\Phi$).}
\emph{Empirically verified} (see \texttt{harmonic\_g2\_factor.py} stabilizer
probe): every antisymmetric generator $U(e_i, e_j) = E_{ij} - E_{ji}$ with
$i, j \in \{1, \ldots, 7\}$ — \emph{all 21 directions inside $\mathrm{so}(7)$} —
preserves $\Phi(\gamma, \gamma, \alpha)$ exactly when $(\alpha, \gamma)$ is in
the standard embedding. In particular $\Delta = E_{17} - E_{71} \in \mathrm{so}(7)$
acts as a $U(1)$ symmetry of the cubic, and the cubic readout
$\Phi = q \cdot a_0$ is a conserved Noether charge of this $U(1)_\Delta$
harmonic flow. This is the geometric content of ``constrain to $g_2$ in
basis~$7$'' from the design notes.

\paragraph{Layer 2 ($\mathrm{so}(8)/\mathrm{so}(7)$ as the harmonic basis).}
The seven generators $U(e_0, e_j)$ for $j = 1, \ldots, 7$ — i.e.\ the
$\mathrm{so}(8)$ rotations that mix the triality-fixed real axis $e_0$ with the
imaginary axes — \emph{break} the $\mathrm{so}(7)$ stabilizer. For
$j \in \{1, 2, 3\}$ (the associative-block boosts), the breaking takes a
closed sinusoidal form:

\begin{theorem}[Harmonic factor extraction]\label{thm:harmonic-extract}
With $\alpha = (a_0, a_1, a_2, a_3, 0, 0, 0, 0)$ in the standard embedding, let
$R_{0j}(\theta) = \exp(\theta\, U(e_0, e_j))$. Then for $j \in \{1, 2, 3\}$
\begin{equation}\label{eq:sinusoid}
\Phi\bigl(R_{0j}(\theta)\gamma,\, R_{0j}(\theta)\gamma,\, R_{0j}(\theta)\alpha\bigr)
 \;=\; q \cdot \bigl(a_0 \cos\theta - a_j \sin\theta\bigr),
\end{equation}
a pure sinusoid with amplitude $A_j := q \cdot \sqrt{a_0^2 + a_j^2}$ and
phase $\arctan(a_j / a_0)$. Consequently
\begin{equation}\label{eq:harmonic-extract}
q \cdot N \;=\; A_1^2 + A_2^2 + A_3^2 \;-\; 2\,\mathrm{DC}^2,
\qquad \mathrm{DC} := \Phi(\gamma, \gamma, \alpha) = q \cdot a_0,
\end{equation}
so $q$ is recovered exactly from four cubic readouts.
\end{theorem}

\begin{proof}[Sketch]
For $j \in \{1, 2, 3\}$, $\gamma$ has no component on $e_0$ or $e_j$, so
$R_{0j}(\theta)\gamma = \gamma$. The rotation acts only on $\alpha$:
$\alpha'_0 = a_0 \cos\theta - a_j \sin\theta$ and $\alpha'_j = a_0 \sin\theta + a_j \cos\theta$;
the other components are fixed. By Theorem~\ref{thm:g2-cubic} (applied with
the rotated $\alpha'$ still in $\mathbb{H}$),
\[
\Phi(\gamma, \gamma, \alpha') = q \cdot \alpha'_0 = q \cdot (a_0 \cos\theta - a_j \sin\theta).
\]
Summing the squared amplitudes:
$A_1^2 + A_2^2 + A_3^2 = q^2(3 a_0^2 + a_1^2 + a_2^2 + a_3^2) = q^2(a_0^2 + a_1^2 + a_2^2 + a_3^2) + 2 q^2 a_0^2 = q^2 p + 2\,\mathrm{DC}^2$, giving \eqref{eq:harmonic-extract}.
\end{proof}

\paragraph{Verification.} The closed-form predictions for $A_j$ match numerical
amplitudes (sweeping $\theta$ over $[0, 2\pi)$) to machine precision; the
extraction \eqref{eq:harmonic-extract} returns the correct $q$ on every
semiprime in our test corpus, including the $a_0 = 0$ cases (where $\mathrm{DC} = 0$
and the extraction reduces to $q \cdot N = \sum_j A_j^2$). See
\texttt{tests/test\_factorization.py::test\_harmonic\_factor\_extract\_closed\_form}.

\paragraph{Scope.}
Like Theorem~\ref{thm:g2-cubic}, this is a structural identity, not a polynomial
factoring algorithm: the readouts $A_j$ and $\mathrm{DC}$ require constructing
$\alpha, \gamma$. The result \emph{is} the precise sense in which the user
description ``harmonic series as $U(1)/\Delta$'' refers to a real geometric
basis: the seven $U(e_0, e_j)$ rotations are the Shor analog of the QFT basis,
and the cubic readout decomposes into sinusoidal harmonics whose amplitudes
encode the Lagrange decomposition of $p$.

\section{What remains outside the formal proof}

\begin{enumerate}[label=\textbf{G\arabic*.}]
\item \textbf{Search completeness:} no Lean theorem shows that either driver
      encounters a divisor for every composite $n$ under stated caps.
\item \textbf{Period $\to$ divisor:} the Shor-style bridge from Fourier/period peaks
      to $\gcd$ factors is not closed inside these files.
\item \textbf{Floating-point:} adequacy of real drift budgets (\texttt{FracDriftAdequateResources})
      is a template; exact Python \texttt{mpmath} semantics are not embedded.
\item \textbf{Primality output:} probable-prime tests in scripts are not theorems.
\item \textbf{$G_2$-cubic readout (Theorem~\ref{thm:g2-cubic}):} the Cayley--Dickson
      identity and the closed forms~\eqref{eq:g2-cubic}--\eqref{eq:comm-norm} are
      verified numerically in \texttt{test\_factorization.py} but not yet expressed
      in Lean. The natural module hosting this would be
      \texttt{Hqiv.Geometry.G2CubicFactorReadout}, parallel to the existing
      \texttt{MonolithicGeometricFactorizer}.
\item \textbf{Harmonic extraction (Theorem~\ref{thm:harmonic-extract}):} the
      sinusoidal closed form~\eqref{eq:sinusoid} and the four-readout
      extraction \eqref{eq:harmonic-extract} are verified numerically against
      direct $\theta$-sweeps but not yet formalized. The
      $\mathrm{so}(7)$-stabilizer claim and the $\mathrm{so}(8)/\mathrm{so}(7)$
      harmonic basis identification are likewise empirical at present.
\end{enumerate}

\section*{Acknowledgments}

Formalization standards and API conventions follow Mathlib~4 \cite{mathlib}.
Statements in this note are intended to track the Lean modules
\texttt{Hqiv.Geometry.FactorDivisibilityBridge},
\texttt{Hqiv.Geometry.HQIVOSHIntegratedFactorDriver},
\texttt{Hqiv.Geometry.MonolithicGeometricFactorizer},
\texttt{Hqiv.Geometry.FactorSearchCostModel}, and
\texttt{Hqiv.QuantumComputing.OSHoracle}. When in doubt, the Lean source is authoritative.

\begin{thebibliography}{9}
\bibitem{mathlib}
The Mathlib community, \emph{Mathlib4}, \url{https://github.com/leanprover-community/mathlib4}.
\end{thebibliography}

\end{document}
